\documentclass[11pt]{article}
% ======================================================================
%  examstyle.tex — EPFL-style exam layout for the weekly Time Series exams
%  Shared by every Exam_Week_NN(.tex / _Solutions.tex). Compiles with pdflatex.
% ======================================================================
\usepackage[utf8]{inputenc}
\usepackage[T1]{fontenc}
\usepackage[english]{babel}
\usepackage[a4paper,margin=2.4cm]{geometry}
\usepackage{amsmath,amssymb,amsthm,mathtools}
\usepackage{bm}
\usepackage{enumitem}
\usepackage{array}

\setlength{\parindent}{0pt}
\setlength{\parskip}{0.5em}
\emergencystretch=3em
\allowdisplaybreaks[1]

% ---------- math macros (same names as the rest of the project) ----------
\newcommand{\E}{\mathbb{E}}
\DeclareMathOperator{\Var}{Var}
\DeclareMathOperator{\Cov}{Cov}
\DeclareMathOperator{\Corr}{Corr}
\DeclareMathOperator*{\argmin}{arg\,min}
\DeclareMathOperator{\MSE}{MSE}
\DeclareMathOperator{\trace}{tr}
\newcommand{\Reals}{\mathbb{R}}
\newcommand{\Ints}{\mathbb{Z}}
\newcommand{\Comp}{\mathbb{C}}
\newcommand{\Nat}{\mathbb{N}}
\newcommand{\Prob}{\mathbb{P}}
\newcommand{\Tr}{^{\mathsf{T}}}
\newcommand{\conj}[1]{\overline{#1}}
\newcommand{\abs}[1]{\left\lvert #1 \right\rvert}
\newcommand{\norm}[1]{\left\lVert #1 \right\rVert}
\newcommand{\ind}{\mathbf{1}}
\newcommand{\eps}{\varepsilon}
\newcommand{\diff}{\nabla}
\newcommand{\acvs}{\gamma}
\newcommand{\acf}{\rho}
\newcommand{\pacf}{\alpha}
\newcommand{\sdf}{\mathcal{S}}
\newcommand{\Bsh}{B}
\newcommand{\iid}{\overset{\text{iid}}{\sim}}

\setlist[enumerate]{leftmargin=2em,topsep=2pt,itemsep=4pt}
\setlist[itemize]{leftmargin=1.6em,topsep=2pt,itemsep=2pt}

\newcounter{exo}

% ---------- exam cover header :  \examheader{exam title}{duration} ----------
\newcommand{\examheader}[2]{%
  \thispagestyle{empty}
  \begin{center}
    {\Large\bfseries Time Series}\\[3pt]
    Spring Semester 2026, EPFL\\
    \'Ecole Polytechnique F\'ed\'erale de Lausanne\\
    Prof.\ Dr.\ Sofia C.\ Olhede\\[7pt]
    \hrule height 0.8pt
    \vspace{9pt}
    {\large\bfseries Practice Exam --- #1}\\[4pt]
    {\normalsize Duration: #2 \quad$\cdot$\quad No calculator \quad$\cdot$\quad Answer on paper}\\
    \vspace{8pt}
    \hrule height 0.8pt
  \end{center}
  \vspace{14pt}
  \begin{tabular}{@{}p{8.2cm}@{}p{8cm}@{}}
    Name:\enspace\hrulefill & Surname:\enspace\hrulefill
  \end{tabular}\\[14pt]
  SCIPER (Student Card Number):\enspace\hrulefill
  \vspace{16pt}
}

% ---------- points table (fixed: 4 problems) ----------
\newcommand{\pointstable}{%
  \begin{center}
  \renewcommand{\arraystretch}{1.5}
  \begin{tabular}{|p{5cm}|p{4cm}|}
    \hline
    \textbf{Exercise} & \textbf{Points}\\\hline
    1 & \\\hline
    2 & \\\hline
    3 & \\\hline
    4 & \\\hline
    \textbf{Total} & \\\hline
  \end{tabular}
  \end{center}
}

% ---------- problem heading (exam: one problem per page) ----------
\newcommand{\problem}[1]{%
  \clearpage
  \stepcounter{exo}%
  \begin{center}\textbf{\large Problem \theexo\ (#1 points)}\end{center}
  \vspace{4pt}
}
% ---------- answer space: fill the statement page + one blank answer page ----------
% Put \answerpage at the END of each problem so every exercise gets its own page(s).
\newcommand{\answerpage}{\par\vfill\newpage\null\newpage}

% ---------- solutions document ----------
\newcommand{\solheader}[1]{%
  \begin{center}
    {\Large\bfseries Time Series --- Solutions}\\[3pt]
    {\large Practice Exam --- #1}\\[2pt]
    EPFL, MATH-342 \quad$\cdot$\quad Prof.\ S.\ C.\ Olhede\\[5pt]
    \hrule height 0.8pt
  \end{center}
  \vspace{10pt}
}
\newcommand{\solproblem}[1]{%
  \bigskip
  \stepcounter{exo}%
  {\large\bfseries Problem \theexo\ (#1 points)}\par\nobreak\vspace{2pt}
}
% Rappel de l'énoncé avant la solution (dans les corrigés)
\newcommand{\statementstart}{\par\vspace{1pt}{\bfseries\itshape Statement.}\par\nobreak\begingroup\small\itshape}
\newcommand{\statementend}{\par\endgroup\vspace{5pt}\hrule\vspace{6pt}{\bfseries Solution.}\par\nobreak\vspace{2pt}}

\begin{document}

\solheader{Week 3: Backshift, Wold Decomposition, Causality \& Invertibility}

% ---------------------------------------------------------------
\solproblem{5}
\statementstart
\textbf{(a)} Define the \emph{backshift operator} $\Bsh$ and, for an $\mathrm{ARMA}(p,q)$ process written $\Phi(\Bsh)X_t=\Theta(\Bsh)\eps_t$, write down the characteristic polynomials $\Phi(z)$ and $\Theta(z)$. State the \emph{root criteria} (in terms of the unit circle) for (i) stationarity/causality and (ii) invertibility. \hfill(2 pts)

\textbf{(b)} Consider the $\mathrm{AR}(2)$ process
\[
X_t=\tfrac23 X_{t-1}-\tfrac16 X_{t-2}+\eps_t .
\]
Write its AR polynomial $\Phi(z)$, find the roots, and decide whether the process is stationary. \hfill(2 pts)

\textbf{(c)} Is the process of part (b) invertible? Justify in one sentence (no computation needed). \hfill(1 pt)
\statementend

\textbf{(a)} The \emph{backshift operator} $\Bsh$ shifts a series one step into the past, $\Bsh X_t=X_{t-1}$, and iterating $\Bsh^k X_t=X_{t-k}$ (with $\Bsh^0=I$). For an $\mathrm{ARMA}(p,q)$ written $\Phi(\Bsh)X_t=\Theta(\Bsh)\eps_t$ the characteristic polynomials are
\[
\Phi(z)=1-\phi_1 z-\dots-\phi_p z^{p},\qquad
\Theta(z)=1+\theta_1 z+\dots+\theta_q z^{q}
\]
(coefficients matching the chosen sign convention). The root criteria are:
\begin{itemize}
  \item \textbf{Stationarity / causality:} all roots of $\Phi(z)$ lie \emph{outside} the unit circle, $\abs{z}>1$;
  \item \textbf{Invertibility:} all roots of $\Theta(z)$ lie \emph{outside} the unit circle, $\abs{z}>1$.
\end{itemize}

\textbf{(b)} In backshift form $\big(1-\tfrac23\Bsh+\tfrac16\Bsh^2\big)X_t=\eps_t$, so
\[
\Phi(z)=1-\tfrac23 z+\tfrac16 z^2 .
\]
Solve $\Phi(z)=0$. Multiplying by $6$ gives $z^2-4z+6=0$, hence
\[
z=\frac{4\pm\sqrt{16-24}}{2}=2\pm\sqrt{-2}=2\pm\sqrt2\,i,
\qquad
\abs{z}=\sqrt{2^2+(\sqrt2)^2}=\sqrt{4+2}=\sqrt6\approx2.449 .
\]
Both roots have modulus $\sqrt6>1$, i.e.\ lie outside the unit circle, so the process is \textbf{stationary} (causal). (One may also check the $\mathrm{AR}(2)$ triangle with $\phi_1=\tfrac23,\ \phi_2=-\tfrac16$: $\phi_1+\phi_2=\tfrac12<1$, $\phi_2-\phi_1=-\tfrac56<1$, $\abs{\phi_2}=\tfrac16<1$.)

\textbf{(c)} \textbf{Yes, invertible.} It is a pure $\mathrm{AR}$ process ($\Theta\equiv1$, a finite polynomial with no roots to test), so $\eps_t=\Phi(\Bsh)X_t$ is already a finite expression in present and past $X$: every $\mathrm{AR}$ is automatically invertible.

% ---------------------------------------------------------------
\solproblem{5}
\statementstart
For the $\mathrm{ARMA}(2,1)$ process
\[
(1-0.7\Bsh+0.1\Bsh^2)\,X_t=(1+0.5\Bsh)\,\eps_t,
\]

\textbf{(a)} factor the AR polynomial $\Phi(z)=1-0.7z+0.1z^2$, locate its roots, and confirm that the process is causal. \hfill(2 pts)

\textbf{(b)} Using a \emph{partial-fraction} expansion of the transfer function $G(z)=\Theta(z)/\Phi(z)$, find a closed form for the $\mathrm{MA}(\infty)$ weights $\psi_j$ in $X_t=\sum_{j\ge0}\psi_j\eps_{t-j}$, and report $\psi_0,\psi_1,\psi_2$. \hfill(3 pts)
\statementend

\textbf{(a)} $\Phi(z)=1-0.7z+0.1z^2=0.1\,(z^2-7z+10)=0.1\,(z-2)(z-5)$, so the AR roots are $z=2$ and $z=5$. Both satisfy $\abs{z}>1$, hence all roots of $\Phi$ are outside the unit circle and the process is \textbf{causal} (stationary). Also $\Theta(z)=1+0.5z$ has root $z=-2$, $\abs{z}=2>1$, so it is invertible as well (not asked, but it guarantees a well-defined $G$).

\textbf{(b)} Write $X_t=G(\Bsh)\eps_t$ with
\[
G(z)=\frac{1+0.5z}{0.1\,(z-2)(z-5)} .
\]
Partial-fraction the rational factor:
\[
\frac{1}{(z-2)(z-5)}=\frac{A}{z-2}+\frac{B}{z-5},\qquad
A=\frac{1}{2-5}=-\frac13,\quad B=\frac{1}{5-2}=\frac13 .
\]
Convert each pole to a power series convergent on (and inside) the unit disk, using $\abs{z}<2<5$:
\[
\begin{aligned}
\frac{-1}{z-2}&=\frac{1}{2}\,\frac{1}{1-z/2}=\frac12\sum_{j\ge0}\Big(\tfrac12\Big)^{j}z^{j}
=\sum_{j\ge0}\Big(\tfrac12\Big)^{j+1}z^{j},\\
\frac{1}{z-5}&=-\frac{1}{5}\,\frac{1}{1-z/5}=-\frac15\sum_{j\ge0}\Big(\tfrac15\Big)^{j}z^{j}
=-\sum_{j\ge0}\Big(\tfrac15\Big)^{j+1}z^{j}.
\end{aligned}
\]
Hence, collecting $\dfrac{1}{(z-2)(z-5)}=\dfrac{A(-1)}{2-z}\cdots$ explicitly,
\[
\frac{1}{(z-2)(z-5)}
=\Big(-\tfrac13\Big)\frac{1}{z-2}+\tfrac13\,\frac{1}{z-5}
=\sum_{j\ge0}\Big[\tfrac13\Big(\tfrac12\Big)^{j+1}-\tfrac13\Big(\tfrac15\Big)^{j+1}\Big]z^{j}.
\]
Now multiply by $\dfrac{1+0.5z}{0.1}=10\,(1+\tfrac12 z)$. Write $H(z)=\sum_{j\ge0}h_j z^j$ with
\[
h_j=10\Big[\tfrac13\Big(\tfrac12\Big)^{j+1}-\tfrac13\Big(\tfrac15\Big)^{j+1}\Big]
=\frac{5}{3}\Big(\tfrac12\Big)^{j}-\frac{2}{3}\Big(\tfrac15\Big)^{j},
\]
so that $G(z)=(1+\tfrac12 z)H(z)$ gives $\psi_0=h_0$ and $\psi_j=h_j+\tfrac12 h_{j-1}$ for $j\ge1$. Computing the coefficients of $(1/2)^j$ and $(1/5)^j$:
\[
\begin{aligned}
\text{coeff of }(\tfrac12)^j:&\quad \tfrac53+\tfrac12\cdot\tfrac53\cdot 2=\tfrac53+\tfrac53=\tfrac{10}{3},\\
\text{coeff of }(\tfrac15)^j:&\quad -\tfrac23+\tfrac12\cdot\big(-\tfrac23\big)\cdot 5=-\tfrac23-\tfrac53=-\tfrac73,
\end{aligned}
\]
(using $(\tfrac12)^{j-1}=2(\tfrac12)^j$ and $(\tfrac15)^{j-1}=5(\tfrac15)^j$). Therefore
\[
\boxed{\;\psi_j=\frac{10}{3}\Big(\tfrac12\Big)^{j}-\frac{7}{3}\Big(\tfrac15\Big)^{j},\qquad j\ge0\;}
\]
which already gives $\psi_0=\tfrac{10}{3}-\tfrac73=1$. In particular
\[
\psi_0=1,\qquad
\psi_1=\tfrac{10}{3}\cdot\tfrac12-\tfrac73\cdot\tfrac15=\tfrac53-\tfrac{7}{15}=\tfrac{6}{5}=1.2,
\qquad
\psi_2=\tfrac{10}{3}\cdot\tfrac14-\tfrac73\cdot\tfrac1{25}=\tfrac56-\tfrac{7}{75}=\tfrac{37}{50}=0.74 .
\]
\textit{Check (coefficient matching).} From $(1-0.7z+0.1z^2)\sum_j\psi_j z^j=1+0.5z$: $\psi_0=1$; $\psi_1-0.7\psi_0=0.5\Rightarrow\psi_1=1.2$; and $\psi_j=0.7\psi_{j-1}-0.1\psi_{j-2}$ for $j\ge2$, giving $\psi_2=0.7(1.2)-0.1(1)=0.74$, matching exactly.

% ---------------------------------------------------------------
\solproblem{5}
\statementstart
Let $\{X_t\}$ be the causal $\mathrm{AR}(2)$ process
\[
X_t=\tfrac56 X_{t-1}-\tfrac16 X_{t-2}+\eps_t .
\]

\textbf{(a)} Show that the autocovariance sequence satisfies the homogeneous recurrence
$\acvs_\tau=\tfrac56\,\acvs_{\tau-1}-\tfrac16\,\acvs_{\tau-2}$ for $\tau\ge1$, and obtain the boundary relation at $\tau=0$ involving $\sigma^2$. \hfill(2 pts)

\textbf{(b)} Solve for $\acvs_0,\acvs_1$ and hence give the closed-form expression of $\acvs_\tau$ for all $\tau\ge0$. \hfill(3 pts)

\emph{Hint: the right-hand side of the model equation only contributes through $\E[\eps_t X_t]=\sigma^2$ (causality).}
\statementend

\textbf{(a)} The model is $X_t-\tfrac56 X_{t-1}+\tfrac16 X_{t-2}=\eps_t$, with $\Phi(z)=1-\tfrac56 z+\tfrac16 z^2=\tfrac16(z-2)(z-3)$ (roots $2,3$ outside the unit circle $\Rightarrow$ causal). Multiply the model by $X_{t-\tau}$ and take expectations. Since the process is causal, $X_{t-\tau}$ is a one-sided combination of $\eps_{t-\tau},\eps_{t-\tau-1},\dots$, so $\E[\eps_t X_{t-\tau}]=0$ for $\tau\ge1$ and $\E[\eps_t X_t]=\psi_0\sigma^2=\sigma^2$. Hence, for $\tau\ge1$,
\[
\acvs_\tau-\tfrac56\acvs_{\tau-1}+\tfrac16\acvs_{\tau-2}=\E[\eps_t X_{t-\tau}]=0
\quad\Longrightarrow\quad
\acvs_\tau=\tfrac56\,\acvs_{\tau-1}-\tfrac16\,\acvs_{\tau-2},
\]
(using $\acvs_{-1}=\acvs_1$ when $\tau=1$). The boundary relation at $\tau=0$ keeps the noise term:
\[
\acvs_0-\tfrac56\acvs_{1}+\tfrac16\acvs_{2}=\E[\eps_t X_t]=\sigma^2 .
\]

\textbf{(b)} The recurrence at $\tau=1,2$ (with $\acvs_{-1}=\acvs_1$) gives the autocorrelations via Yule--Walker:
\[
\acf_1=\frac{\phi_1}{1-\phi_2}=\frac{5/6}{1+1/6}=\frac{5/6}{7/6}=\frac57,
\qquad
\acf_2=\phi_1\acf_1+\phi_2=\tfrac56\cdot\tfrac57-\tfrac16=\tfrac{25}{42}-\tfrac{7}{42}=\tfrac{18}{42}=\tfrac37 .
\]
Plug into the $\tau=0$ relation, $\acvs_0\big(1-\phi_1\acf_1-\phi_2\acf_2\big)=\sigma^2$:
\[
1-\tfrac56\cdot\tfrac57-\big(-\tfrac16\big)\tfrac37=1-\tfrac{25}{42}+\tfrac{3}{42}=1-\tfrac{22}{42}=\tfrac{20}{42}=\tfrac{10}{21},
\qquad
\acvs_0=\frac{21}{10}\sigma^2 .
\]
Therefore $\acvs_1=\acf_1\acvs_0=\tfrac57\cdot\tfrac{21}{10}\sigma^2=\tfrac32\sigma^2$ and $\acvs_2=\acf_2\acvs_0=\tfrac37\cdot\tfrac{21}{10}\sigma^2=\tfrac9{10}\sigma^2$.

The homogeneous recurrence has characteristic equation $\Phi$, whose roots correspond to the reciprocal roots $\tfrac12,\tfrac13$, so the general solution is $\acvs_\tau=c_1\big(\tfrac12\big)^{\tau}+c_2\big(\tfrac13\big)^{\tau}$. Match $\acvs_0=\tfrac{21}{10}\sigma^2$ and $\acvs_1=\tfrac32\sigma^2$:
\[
\begin{aligned}
c_1+c_2&=\tfrac{21}{10}\sigma^2,\\
\tfrac12 c_1+\tfrac13 c_2&=\tfrac32\sigma^2 .
\end{aligned}
\]
Subtracting $\tfrac13\times$(first) from the second, $\tfrac16 c_1=\big(\tfrac32-\tfrac{21}{30}\big)\sigma^2=\tfrac{45-21}{30}\sigma^2=\tfrac{24}{30}\sigma^2$, so $c_1=\tfrac{24}{5}\sigma^2$ and $c_2=\tfrac{21}{10}\sigma^2-\tfrac{24}{5}\sigma^2=-\tfrac{27}{10}\sigma^2$. Hence
\[
\boxed{\;\acvs_\tau=\Big[\tfrac{24}{5}\big(\tfrac12\big)^{\tau}-\tfrac{27}{10}\big(\tfrac13\big)^{\tau}\Big]\sigma^2,\qquad \tau\ge0\;}
\]
and $\acvs_{-\tau}=\acvs_\tau$. (Check: $\tau=0\to\tfrac{48-27}{10}=\tfrac{21}{10}$; $\tau=1\to\tfrac{24}{10}-\tfrac{27}{30}=\tfrac{72-27}{30}=\tfrac{45}{30}=\tfrac32$; $\tau=2\to\tfrac{24}{20}-\tfrac{27}{90}=\tfrac{6}{5}-\tfrac{3}{10}=\tfrac{9}{10}$, all matching.)\\
\textit{Independent check (via $\psi$-weights):} $\psi_j=2\big(\tfrac12\big)^j-\big(\tfrac13\big)^j$ solves $\psi_j=\tfrac56\psi_{j-1}-\tfrac16\psi_{j-2}$ with $\psi_0=1,\psi_1=\tfrac56$, and $\sigma^2\sum_{j\ge0}\psi_j\psi_{j+\tau}$ reproduces $\tfrac{21}{10},\tfrac32,\tfrac9{10}\dots$ exactly.

% ---------------------------------------------------------------
\solproblem{5}
\statementstart
\textbf{(Revision of Week 2 + this week.)} Let $X_t=\eps_t-\theta\eps_{t-1}$ be an $\mathrm{MA}(1)$ with $\theta=\tfrac12$.

\textbf{(a)} Compute the autocovariance sequence $\acvs_\tau$ for all $\tau$ and the autocorrelation $\acf_1$. \hfill(1.5 pts)

\textbf{(b)} Prove that \emph{any} $\mathrm{MA}(1)$ of the form $\eps_t-\theta\eps_{t-1}$ satisfies $\abs{\acf_1}\le\tfrac12$, and state qualitatively how the PACF of an $\mathrm{MA}(1)$ behaves (cut off or tail off?). \hfill(1.5 pts)

\textbf{(c)} Is $\{X_t\}$ invertible? Then exhibit a \emph{different} $\mathrm{MA}(1)$ process $Y_t=\eta_t-\vartheta\eta_{t-1}$ (with an appropriately chosen white-noise variance) that has \emph{exactly the same} autocovariance sequence as $\{X_t\}$, and explain which of the two is the invertible representation. \hfill(2 pts)
\statementend

\textbf{(a)} With $X_t=\eps_t-\tfrac12\eps_{t-1}$, bilinearity and $\Cov(\eps_s,\eps_u)=\sigma^2\ind_{\{s=u\}}$ give
\[
\acvs_\tau=\begin{cases}
(1+\theta^2)\sigma^2=\big(1+\tfrac14\big)\sigma^2=\tfrac54\sigma^2,&\tau=0,\\[2pt]
-\theta\,\sigma^2=-\tfrac12\sigma^2,&\tau=\pm1,\\[2pt]
0,&\abs\tau\ge2,
\end{cases}
\qquad
\acf_1=\frac{-\theta}{1+\theta^2}=\frac{-1/2}{5/4}=-\tfrac25 .
\]

\textbf{(b)} For a general $\mathrm{MA}(1)$, $\acf_1=\dfrac{-\theta}{1+\theta^2}$, so
\[
\abs{\acf_1}=\frac{\abs\theta}{1+\theta^2}\le\frac12
\quad\Longleftrightarrow\quad
2\abs\theta\le 1+\theta^2
\quad\Longleftrightarrow\quad
0\le\theta^2-2\abs\theta+1=(\abs\theta-1)^2,
\]
which always holds (equality iff $\abs\theta=1$). Equivalently, by AM--GM $\abs\theta\le\tfrac12(1+\theta^2)$. So $\abs{\acf_1}\le\tfrac12$.

\emph{PACF behaviour:} for an $\mathrm{MA}(1)$ the \textbf{PACF tails off} (decays geometrically in magnitude with alternating sign, never cutting to exactly $0$), while it is the ACF that cuts off after lag $1$. (This is the diagnostic mirror of the $\mathrm{AR}(1)$, whose PACF cuts off after lag $1$ and ACF tails off.)

\textbf{(c)} \textbf{Invertible:} $\Theta(z)=1-\tfrac12 z$ has root $z=1/\theta=2$ with $\abs{2}>1$ (outside the unit circle), so $\{X_t\}$ is invertible.

A different $\mathrm{MA}(1)$ with the same ACVS is the \emph{reciprocal-root twin}: take coefficient $\vartheta=1/\theta=2$ and white-noise variance $\sigma_\eta^2=\theta^2\sigma^2=\tfrac14\sigma^2$, i.e.
\[
Y_t=\eta_t-2\eta_{t-1},\qquad \Var(\eta_t)=\tfrac14\sigma^2 .
\]
Its autocovariances are
\[
\acvs^Y_0=(1+\vartheta^2)\sigma_\eta^2=(1+4)\cdot\tfrac14\sigma^2=\tfrac54\sigma^2,
\qquad
\acvs^Y_{\pm1}=-\vartheta\,\sigma_\eta^2=-2\cdot\tfrac14\sigma^2=-\tfrac12\sigma^2,
\]
and $\acvs^Y_\tau=0$ for $\abs\tau\ge2$ --- \emph{identical} to $\{X_t\}$ in part (a). So the two are indistinguishable from second-order structure alone (MA non-identifiability). The \textbf{invertible} representation is $\{X_t\}$, whose MA root $z=2$ lies \emph{outside} the unit circle; the twin $\{Y_t\}$ has root $z=\vartheta^{-1}=\tfrac12$ \emph{inside} the circle and is therefore \emph{not} invertible. Imposing invertibility selects $\{X_t\}$ as the unique model.

\end{document}
